\documentclass[12pt,a4paper]{article}

\usepackage{amsmath,amssymb,amsthm}
\usepackage{enumerate}
\usepackage{enumitem}
\usepackage{hyperref}
\usepackage{geometry}
\usepackage{doi}
\geometry{margin=2.5cm}

\hypersetup{
  hidelinks,
  unicode,
  pdftitle={Quadratic Completion of Admissible Spectral Pairs via Binary-Icosahedral Representation},
  pdfauthor={Jérôme Beau}
}
\setlength{\emergencystretch}{2em}

% Theorem environments
\newtheorem{theorem}{Theorem}[section]
\newtheorem{proposition}[theorem]{Proposition}
\newtheorem{lemma}[theorem]{Lemma}
\newtheorem{corollary}[theorem]{Corollary}
\newtheorem{conjecture}[theorem]{Conjecture}
\newtheorem{hypothesis}[theorem]{Hypothesis}
\theoremstyle{definition}
\newtheorem{definition}[theorem]{Definition}
\newtheorem{criterion}[theorem]{Criterion}
\theoremstyle{remark}
\newtheorem{remark}[theorem]{Remark}

\title{Quadratic Completion of Admissible Spectral Pairs\\
via Binary-Icosahedral Representation\\[0.5em]
\large O26 --- Spectral Admissibility Sub-programme}

\author{J\'er\^ome Beau\\
\small Independent Researcher}

\date{2026\\[2pt]\small Version 2.1}

\begin{document}

  \maketitle

  \begin{abstract}
    We investigate the structural nature of the pair observable
    $\sigma_{\mathrm{pair}}(n) = \sigma_c(n)\,\sigma_{q-c}(n)$ arising in the
    Weil-block analysis of the Heisenberg Cayley graphs
    $\mathrm{Heis}_3(\mathbb{Z}/q\mathbb{Z})$.
    Building on the canonical construction of conjugate pairs under the parity
    involution $c \leftrightarrow q-c$, the candidate fibre structure of the pair
    sector (a structurally motivated hypothesis; O18 states the
    fibre identification as an open problem), and the normalisation invariance
    established in O17--O19,
    we construct an explicit dictionary between conjugate Weil blocks and rank-one
    matrix coefficients in a representation space.

    We show that, in the pre-saturation regime, $\sigma_{\mathrm{pair}}(n)$ has the
    same growth exponent as the Hilbert--Schmidt norm of an associated matrix
    trajectory, establishing a Level~I identification (proved).
    We then formulate a hierarchy of stronger identifications:
    a quotient identification modulo normalisation (Level~II), and a canonical
    representation-theoretic identification (Level~III), which is a theorem
    conditional on a single structural hypothesis (the admissible embedding
    $\Phi_{q,\rho}$) in an isotypic sector of the binary icosahedral group $2I$
    along the admissibility thread $Q_8 \subset 2I \subset \mathrm{SU}(2)$.

    Under the involution-equivariance condition of that embedding, every Level~III outer
    product is a symmetric square $uu^{\mathsf T}$, so the trajectory spans at most
    $\mathrm{Sym}^2(V_\rho)$, of dimension $3$ when $\dim_{\mathbb{C}} V_\rho = 2$: the
    rank-four target in $\mathrm{End}(V_\rho)$ is excluded under that contract, without
    refuting other embeddings or constructions.
    The conjugate-pair covariance computable from O25 data is formed in the measured carrier
    $H_{\mathrm{eff}}$ rather than in $\mathrm{End}(V_\rho)$; it yields finite-data diagnostics
    compatible with a supplied adjoint carrier and does not identify a representation sector.
  \end{abstract}

  \paragraph{Keywords.}
    spectral admissibility; pair observable; Hilbert–Schmidt norm; Weil representation; conjugate
    pairs; representation theory; non-injective projection

    \clearpage
    \tableofcontents

    \clearpage
  \section{Background and motivation}
  \label{sec:background}

  The spectral admissibility programme (O1--O25) studies the emergence of observable
  scaling laws from the non-injective projection structure of Cosmochrony~\cite{Cosmochrony},
  in which the static substrate~$\chi$ is projected onto observable spacetime via~$\Pi$.
  A central object is the Gram--Schmidt redundancy observable
  \begin{equation}
    \sigma_c(n) = \frac{\delta r_n}{|S_n|},
  \end{equation}
  defined on BFS shells $S_n$ of Cayley graphs associated with
  $\mathrm{Heis}_3(\mathbb{Z}/q\mathbb{Z})$.

  Early analyses based on scalar observables led to systematic discrepancies in the
  extracted exponent~$\delta$ (O11--O14~\cite{Beau2026a15,Beau2026a16,Beau2026a17,Beau2026a18}).
  This mismatch was resolved in O16~\cite{Beau2026a20} by identifying the correct observable as
  the \emph{conjugate pair observable}
  \begin{equation}
    \sigma^{\mathrm{can}}_{\mathrm{pair}}(n) = \sigma_c(n)\,\sigma_{q-c}(n),
  \end{equation}
  which doubles the exponent and yields values in the admissible window
  $\delta_{\mathrm{pair}} \in [7.4, 10.6]$.

  The pairing rests on the conjugation identity
  $\rho_{q-c} = \overline{\rho_c}$ (an O17 theorem, recalled in O18 Theorem~3.1).
  O18~\cite{Beau2026a22} establishes Born--Infeld parity as a covariance of the response family
  of the substrate action; the reading of the parity orbit $\{c, q-c\}$ as a fibre
  of $\Pi$ is stated there as an open problem, so the pair observable keeps the
  status of a structurally motivated hypothesis.
  O17--O19~\cite{Beau2026a21,Beau2026a22,Beau2026a23} establish that the apparent
  asymmetry between $c$ and $q-c$ is a normalisation artefact, not a structural feature.

  The present note addresses the following question:
  \begin{quote}
    Is the bilinear form $\sigma^{\mathrm{can}}_{\mathrm{pair}}(n)$ an intrinsic
    quadratic object, or merely a product enforced by symmetry?
  \end{quote}
  We show that $\sigma^{\mathrm{can}}_{\mathrm{pair}}$ admits a natural interpretation
  as the pullback of a Hermitian quadratic form on a representation space, and formulate
  precise criteria under which this interpretation becomes canonical.

  \section{The admissible $2I$ thread and the 600-cell}
  \label{sec:2I}

  The admissibility analysis of~\cite{Beau2026a1,Beau2026a3} classifies the finite
  subgroups of $\mathrm{SU}(2)$ compatible with bounded-flux admissibility constraints,
  identifying the chain
  \begin{equation}
    Q_8 \subset 2I \subset \mathrm{SU}(2),
  \end{equation}
  where $2I$ is the binary icosahedral group of order $120$.
  Among the admissible families, the binary polyhedral groups $2O$ and $2I$ are
  distinguished by the isotropy of their second-moment tensor~\cite{Beau2026a3},
  making $2I$ the maximal admissible finite subgroup of $\mathrm{SU}(2)$.

  The group $2I$ admits irreducible representations of dimensions
  \begin{equation}
    d_\rho \in \{1, 2, 2, 3, 3, 4, 4, 5, 6\},
  \end{equation}
  and its matrix coefficient spaces $\mathrm{End}(V_\rho)$ have dimensions $d_\rho^2$,
  forming the canonical Peter--Weyl decomposition of $L^2(2I)$.

  Geometrically, $2I$ is the rotational symmetry group of the regular 600-cell in
  $\mathbb{R}^4$, whose $120$ vertices correspond to the elements of $2I$.
  The Laplacian spectrum of this structure can be organised representation-theoretically,
  with multiplicity patterns consistent with the Peter--Weyl dimensions $d_\rho^2$ of
  the irreps of $2I$~\cite{Beau2026a1}.
  This is not a coincidence imported from external geometry: the same representation-theoretic
  hierarchy governs both the spectral decomposition of $2I$ and the organisation
  of admissible degrees of freedom along the thread $Q_8 \subset 2I \subset \mathrm{SU}(2)$.

  The question addressed in this note is whether the conjugate-pair observable
  $\sigma^{\mathrm{can}}_{\mathrm{pair}}$ selects one of these isotypic sectors and
  realises its canonical quadratic form through the Gram--Schmidt dynamics.

  \section{The dictionary: conjugate Weil blocks and isotypic sectors of $2I$}
  \label{sec:dictionary}

  This section establishes the explicit correspondence between the pair observable
  $\sigma^{\mathrm{can}}_{\mathrm{pair}}$ on conjugate Weil blocks $\{c, q-c\}$
  and a natural Hermitian quadratic form on an isotypic sector of $2I$.
  The dictionary has three layers: the Weil side (§\ref{sec:dict-weil}),
  the $2I$ side (§\ref{sec:dict-2I}), and the identification map (§\ref{sec:dict-map}).

  \subsection{The Weil side: structure of $\sigma^{\mathrm{can}}_{\mathrm{pair}}$}
    \label{sec:dict-weil}

    Fix a prime $q$ and a generic block $c_{\mathrm{block}} = (c_1, c_2, c_3)$
    with $c_i \neq 0$ and $c_1 + c_2 + c_3 \not\equiv 0 \pmod{q}$.
    The Gram--Schmidt redundancy at BFS shell $n$ is
    \begin{equation}
      \sigma_c(n) = \frac{\delta r_n}{|S_n|},
    \end{equation}
    where $\delta r_n$ counts the number of Weil fingerprint vectors in shell $S_n$
    linearly independent of the span built over shells $0, \ldots, n-1$~\cite{Beau2026a16}.
    The canonical pair observable is
    \begin{equation}
      \sigma^{\mathrm{can}}_{\mathrm{pair}}(n) = \sigma_c(n)\,\sigma_{q-c}(n).
    \end{equation}

    The structural basis of this definition rests on two results from O17--O18.
    First, the conjugation identity $\rho_{q-c} = \overline{\rho_c}$ together with
    the norm-invariance of Gram--Schmidt orthogonalisation implies $\delta_{q-c} = \delta_c$
    for all admissible~$c$~\cite{Beau2026a21}.
    Second, the proportionality $\sigma_{q-c}(n) = r(c,q)\,\sigma_c(n)$, with
    $r(c,q) \in \{1,2,3,4\}$ independent of $n$, holds to machine precision
    ($|\sigma_{q-c}(n)/\sigma_c(n) - r(c,q)| < 10^{-14}$ for all tested pairs
    and shells)~\cite{Beau2026a20}.
    Critically, O17~\cite{Beau2026a21} demonstrates that $r(c,q)$ is a normalisation
    artefact of the pipeline determined by the initial supports $(b_1, b_2)$ of each
    block sample, not an intrinsic invariant of the representation.
    The structurally robust quantity is therefore $\delta_c = \delta_{q-c}$, and the
    canonical form of the pair observable removes the artefact by working with
    the product $\sigma_c(n)\,\sigma_{q-c}(n)$ rather than each factor in isolation.

    \begin{remark}[Gauge-like nature of $r(c,q)$]
  The factor $r(c,q)$ plays the role of a gauge degree of freedom in the pipeline:
  it depends on sampling initial conditions, not on the representation.
  The matrix $M_n$ defined in §\ref{sec:dict-map} must be constructed so that
  $\|\widetilde{M}_n\|_{\mathrm{HS}}^2$ is independent of this gauge choice,
  i.e.\ it absorbs $r(c,q)$ into a normalisation convention rather than treating
  it as structural data.
    \end{remark}

  \subsection{The $2I$ side: isotypic sectors and the Peter--Weyl quadratic form}
    \label{sec:dict-2I}

    For each irrep $\rho$ of $2I$ with representation space $V_\rho$ of dimension
    $d_\rho$, the Peter--Weyl theorem provides a canonical space of matrix coefficients:
    \begin{equation}
      \mathcal{H}_\rho
      = \{f : 2I \to \mathbb{C} \mid f(g) = \langle u, \rho(g) v \rangle,\;
      u, v \in V_\rho\}
      \simeq V_\rho \otimes \overline{V_\rho}
      \simeq \mathrm{End}(V_\rho),
    \end{equation}
    which has dimension $d_\rho^2$.
    This space carries a canonical Hermitian inner product, the Hilbert--Schmidt form:
    \begin{equation}
      \langle A, B \rangle_{\mathrm{HS}} = \mathrm{tr}(A^\dagger B),
      \qquad \|A\|_{\mathrm{HS}}^2 = \mathrm{tr}(A^\dagger A),
      \quad A, B \in \mathrm{End}(V_\rho).
      \label{eq:HS-form}
    \end{equation}
    For a rank-one element $A = u \otimes \overline{v}$, this reduces to
    \begin{equation}
      \|u \otimes \overline{v}\|_{\mathrm{HS}}^2 = \|u\|^2 \cdot \|v\|^2.
      \label{eq:HS-rank1}
    \end{equation}
    This is the key structural fact: the Hilbert--Schmidt norm of a rank-one
    endomorphism is the product of the norms of its two factors.

  \subsection{The dictionary map}
    \label{sec:dict-map}

    \paragraph{Construction.}
      For each BFS shell $n$ and each conjugate pair $\{c, q-c\}$, define the
      \emph{un-normalised matrix coefficient}
      \begin{equation}
        \widetilde{M}_n := v_c^{(n)} \otimes \overline{v_{q-c}^{(n)}},
        \label{eq:Mn-unnorm}
      \end{equation}
      where $v_c^{(n)}$ is the vector whose $s$-th component is the residual norm
      $\|r_s^{(n)}\|$ of the $s$-th Weil fingerprint vector in shell $S_n$ after
      projection onto the Gram--Schmidt span built over shells $0, \ldots, n-1$,
      and similarly for $v_{q-c}^{(n)}$.

    \paragraph{Central identity.}
      By~\eqref{eq:HS-rank1},
      \begin{equation}
        \|\widetilde{M}_n\|_{\mathrm{HS}}^2
        = \|v_c^{(n)}\|^2 \cdot \|v_{q-c}^{(n)}\|^2
        \sim \sigma_c(n) \cdot \sigma_{q-c}(n)
        = \sigma^{\mathrm{can}}_{\mathrm{pair}}(n).
        \label{eq:identity}
      \end{equation}

      \begin{remark}[From pipeline quantity to structural norm]
  The proportionality $\|v_c^{(n)}\|^2 \sim \sigma_c(n)$ in~\eqref{eq:identity}
  is not an additional assumption; it follows from the definition of the
  Gram--Schmidt observable and a single normalisation choice.

  With the convention above, $\|v_c^{(n)}\|^2 = \sum_{s \in S_n} \|r_s^{(n)}\|^2$,
  the total squared residual energy contributed by shell $S_n$.
  By construction, $\sigma_c(n) = \delta r_n / |S_n|$, where $\delta r_n$ is the
  rank increment of shell $S_n$.

  Under the canonical normalisation of the fingerprint vectors adopted in
  O12--O19~\cite{Beau2026a16,Beau2026a21,Beau2026a22,Beau2026a23}, each fingerprint
  vector has unit initial norm.
  The proportionality
  \begin{equation}
    \|v_c^{(n)}\|^2 = \kappa \cdot |S_n| \cdot \sigma_c(n) + \varepsilon_n
    \label{eq:norm-sigma-relation}
  \end{equation}
  holds with $\kappa > 0$ a normalisation constant independent of $n$, and
  $\varepsilon_n$ a remainder bounded by the contribution of vectors below the
  rank threshold.
  In the pre-saturation regime (the fitting window $[n_0, n_1]$ used throughout
  O16--O25), $\varepsilon_n / (|S_n| \cdot \sigma_c(n)) \to 0$ as the
  rank-threshold vectors become negligible relative to the total shell size.

  Identity~\eqref{eq:identity} therefore holds in the pre-saturation regime with
  a proportionality constant $\kappa^2 |S_n|^2$ absorbed into the pipeline
  normalisation.
  The power-law exponent $\delta_{\mathrm{pair}}$, extracted from the log-log
  regression of $\sigma^{\mathrm{can}}_{\mathrm{pair}}(n)$, is unaffected by
  this constant factor.
      \end{remark}

    \paragraph{Three verification criteria.}

      \begin{enumerate}[label=(\roman*)]

\item \textbf{Involution compatibility.}
The exchange $c \leftrightarrow q-c$ maps $\widetilde{M}_n$ to its Hermitian adjoint identically:
by~\eqref{eq:Mn-unnorm}, $(v_c^{(n)} \otimes \overline{v_{q-c}^{(n)}})^\dagger
= v_{q-c}^{(n)} \otimes \overline{v_c^{(n)}}$ for any vectors.
The pair observable $\sigma^{\mathrm{can}}_{\mathrm{pair}}(n) =
\sigma_c(n)\,\sigma_{q-c}(n)$ is symmetric under this exchange, consistent
with $\|A\|_{\mathrm{HS}}^2 = \|A^\dagger\|_{\mathrm{HS}}^2$.
Involution compatibility therefore holds by construction.
It does not use the conjugation identity $\rho_{q-c} = \overline{\rho_c}$, whose scalar content is
$\delta_{q-c} = \delta_c$~\cite{Beau2026a21}, and it does not imply an anti-linear relation between
the coordinates of conjugate blocks: that relation is condition~(a) of
Hypothesis~\ref{hyp:embedding} for the embedded vectors, and its analogue on measured coordinates is
tested in O29~\cite{Beau2026a33}, which finds it satisfied only approximately.

\item \textbf{Gauge invariance.}
The observable $\sigma^{\mathrm{can}}_{\mathrm{pair}}(n)$ is invariant under
rescaling $v_c^{(n)} \mapsto \lambda v_c^{(n)}$,
$v_{q-c}^{(n)} \mapsto \lambda^{-1} v_{q-c}^{(n)}$ with $\lambda > 0$,
because $\sigma_c(n)$ scales as $\lambda^2$ and $\sigma_{q-c}(n)$ as
$\lambda^{-2}$, leaving the product unchanged.
Correspondingly, $\|\widetilde{M}_n\|_{\mathrm{HS}}^2$ is invariant under
$(u, v) \mapsto (\lambda u, \lambda^{-1} v)$ in the rank-one element
$u \otimes \overline{v}$.
Gauge invariance holds.

\item \textbf{$\delta_{\mathrm{pair}}$ stability.}
The exponent $\delta_{\mathrm{pair}}$ is extracted from the log-log regression
of $\sigma^{\mathrm{can}}_{\mathrm{pair}}(n)$ over the fitting window.
By~\eqref{eq:identity}, this is the same as the log-log regression of
$\|\widetilde{M}_n\|_{\mathrm{HS}}^2$.
The verticality theorem of O24~\cite{Beau2026a28} establishes that
$\delta_{\mathrm{pair}}$ is stable under fibre perturbations; this translates,
via the dictionary, into the statement that $\|\widetilde{M}_n\|_{\mathrm{HS}}^2$
grows with a well-defined power-law exponent independent of the choice of
conjugate pair.
$\delta_{\mathrm{pair}}$ stability holds.

      \end{enumerate}

  \subsection{Strength of the identification}
  \label{sec:dict-strength}

  Identity~\eqref{eq:identity} establishes the correspondence at \emph{intermediate
  strength} in the sense of §\ref{sec:conjecture}: $\sigma^{\mathrm{can}}_{\mathrm{pair}}(n)$
  is proportional to $\|\widetilde{M}_n\|_{\mathrm{HS}}^2$, with a proportionality
  constant that is pipeline-normalisation-dependent and hence not intrinsic.

  The \emph{strong} identification requires the proportionality constant to be
  canonically fixed by the $2I$-representation theory, i.e.\ that $\widetilde{M}_n$
  takes values in a specific isotypic component $\mathcal{H}_\rho$ for a
  representation $\rho$ selected by the admissibility thread
  $Q_8 \subset 2I \subset \mathrm{SU}(2)$.
  The admissibility hierarchy of~\cite{Beau2026a1} selects the spin-$\frac{1}{2}$
  sector (dimension $d_\rho = 2$, Laplacian eigenvalue $\lambda_{1/2} = 18$) as
  the dominant admissible sector on the $2I$ Cayley graph.
  If $v_c^{(n)} \in V_\rho$ for this sector, then $\widetilde{M}_n \in \mathrm{End}(V_\rho)$
  has $d_\rho^2 = 4$ components.
  This identification is at present a conjecture, formulated precisely in §\ref{sec:conjecture}.

  \section{The conjecture at three levels}
  \label{sec:conjecture}

  Building on the dictionary of §\ref{sec:dictionary}, we formulate a graduated
  identification between the pair observable $\sigma^{\mathrm{can}}_{\mathrm{pair}}(n)$
  and the Hilbert--Schmidt quadratic form on an isotypic sector of $2I$.
  The three levels differ in the strength of the claim: growth equivalence (proved),
  quotient identification (structural), and canonical identification (conjectural).

  \subsection{Level I: Growth equivalence (proved)}
    \label{sec:level1}

    \begin{proposition}[Level~I]
  In the pre-saturation regime $n \in [n_0, n_1]$, the pair observable and the
  Hilbert--Schmidt norm of $\widetilde{M}_n$ have identical power-law growth exponents:
  \begin{equation}
    \delta_{\mathrm{pair}} = \delta_{\mathrm{HS}},
    \label{eq:level1}
  \end{equation}
  where $\delta_{\mathrm{HS}}$ is defined by the log-log regression of
  $\|\widetilde{M}_n\|_{\mathrm{HS}}^2$ over the same fitting window.
    \end{proposition}

    \begin{proof}
      By~\eqref{eq:identity} and~\eqref{eq:norm-sigma-relation},
      $\sigma^{\mathrm{can}}_{\mathrm{pair}}(n)
      = \kappa^2 |S_n|^2 \|\widetilde{M}_n\|_{\mathrm{HS}}^2 + O(\varepsilon_n)$,
      where the prefactor $\kappa^2 |S_n|^2$ is independent of the conjugate pair.
      Taking log-log regression over $[n_0, n_1]$ and noting that $|S_n|$ grows
      as a fixed power of $n$ (determined by the BFS geometry of $G_q$, independent
      of $c$), the regressions of $\log \sigma^{\mathrm{can}}_{\mathrm{pair}}$ and
      $\log \|\widetilde{M}_n\|_{\mathrm{HS}}^2$ differ only in their intercepts.
      Hence $\delta_{\mathrm{pair}} = \delta_{\mathrm{HS}}$.
    \end{proof}

    \begin{remark}
      Level~I is not a conjecture but a direct consequence of the dictionary
      construction of §\ref{sec:dictionary}.
      It is stated separately to clarify what is proved and what remains conjectural
      at Levels~II and~III.
    \end{remark}

  \subsection{Level II: Quotient identification (structural)}
    \label{sec:level2}

    The proportionality factor $\kappa^2 |S_n|^2$ has two sources: the global
    normalisation convention of the fingerprint vectors (absorbed in $\kappa$) and
    the shell-size scaling (absorbed in $|S_n|^2$).
    Both are pipeline conventions rather than representation-theoretic data.
    Define the \emph{normalisation equivalence}
    \begin{equation}
      A \approx B \quad \Longleftrightarrow \quad A = C \cdot B,
    \end{equation}
    where $C > 0$ depends only on pipeline conventions and the BFS geometry of $G_q$,
    but not on the conjugate pair $\{c, q-c\}$ or the shell index $n$.

    \begin{conjecture}[Level~II]
      \label{conj:level2}
      In the normalisation equivalence class defined above,
      \begin{equation}
        \sigma^{\mathrm{can}}_{\mathrm{pair}}(n)
        \approx \|\widetilde{M}_n\|_{\mathrm{HS}}^2.
        \label{eq:level2}
      \end{equation}
      Equivalently, $\sigma^{\mathrm{can}}_{\mathrm{pair}}(n)$ is the pullback of
      the Hilbert--Schmidt quadratic form on $\mathrm{End}(V_\rho)$ under the
      dictionary map~\eqref{eq:Mn-unnorm}, modulo a constant depending only on
      pipeline conventions.
    \end{conjecture}

    Level~II would follow from Level~I together with a proof that the proportionality
    constant $\kappa^2 |S_n|^2$ does not depend on $n$ in a way that could alter
    the pair-level structure.
    The $n$-dependence of $|S_n|^2$ is BFS-geometric and uniform across pairs,
    so the primary open point is the $n$-independence of $\kappa$.

  \subsection{Level III: Conditional theorem}
    \label{sec:level3}

    The strong form requires the proportionality constant to be fixed canonically by
    the representation theory of $2I$.
    We show that this identification is a \emph{theorem}, conditional on a single
    structural hypothesis: the existence of a canonical assignment from admissible
    Weil-block residuals to the spin-$\frac{1}{2}$ representation space of $2I$.

    \begin{hypothesis}[Canonical admissible embedding]
      \label{hyp:embedding}
      Let $\mathcal{V}_q$ denote the space generated by admissible Weil-block
      residual vectors $\{v_c^{(n)}\}$, equipped with the parity involution
      $c \mapsto q-c$, and let $\rho$ be the spin-$\frac{1}{2}$ irreducible
      representation of $2I$ with representation space $V_\rho$ (dimension $d_\rho = 2$).

      There exists a canonical linear map
      \begin{equation}
        \Phi_{q,\rho} : \mathcal{V}_q \longrightarrow V_\rho
        \label{eq:embedding}
      \end{equation}
      satisfying:
      \begin{enumerate}[label=(\alph*)]
    \item \textbf{Involution equivariance.}
    The parity involution on $\mathcal{V}_q$ maps to the canonical
    anti-linear involution on $V_\rho$ induced by Hermitian conjugation
    in $\mathrm{End}(V_\rho)$:
    \begin{equation}
      \Phi_{q,\rho}(v_{q-c}^{(n)}) = \overline{\Phi_{q,\rho}(v_c^{(n)})}.
      \label{eq:hyp-involution}
    \end{equation}
    \item \textbf{Quadratic compatibility.}
    The pullback of the Hilbert--Schmidt form satisfies
    \begin{equation}
      \|\Phi_{q,\rho}(v_c^{(n)})\|^2 = d_\rho^{-1}\,\|v_c^{(n)}\|^2
      \quad \text{for all } n \in [n_0, n_1].
    \end{equation}
    \item \textbf{Naturality.}
    The assignment $\Phi_{q,\rho}$ is canonical with respect to the
    admissible structure: it is independent of the choice of pair
    $\{c, q-c\}$ and of the shell index $n$, and depends only on
    $(q, \rho)$.
      \end{enumerate}
      This assignment can be viewed as a natural morphism from the space of admissible
      residual configurations to an isotypic component of $2I$.
    \end{hypothesis}

    Hypothesis~\ref{hyp:embedding} isolates the only non-constructive step of the
    argument: all subsequent statements are formal consequences of the constructions
    of §\ref{sec:dictionary}.
    It remains open.
    There is no non-zero $\mathfrak{su}(2)$-equivariant map from the adjoint
    representation to the fundamental representation, so an equivariant adjoint-factorising
    realisation of Hypothesis~\ref{hyp:embedding} is excluded (O27~\cite{Beau2026a31},
    Theorem~4.1(iv)).
    This does not exclude every quotient-factorising linear lift: the condition
    $\ker \Pi_{\mathrm{adm}} \subseteq \ker \Phi$ gives unique linear factorisation and no
    equivariance.
    Hypothesis~\ref{hyp:embedding} therefore remains open in that wider class.

    \begin{theorem}[Level~III: canonical quadratic realisation]
      \label{thm:level3}
      Assume Hypothesis~\ref{hyp:embedding}.
      Then for every admissible conjugate pair $\{c, q-c\}$ and every
      $n \in [n_0, n_1]$,
      \begin{equation}
        \sigma^{\mathrm{can}}_{\mathrm{pair}}(n)
        = C_n \cdot
        \bigl\|\Phi_{q,\rho}(v_c^{(n)}) \otimes
        \overline{\Phi_{q,\rho}(v_{q-c}^{(n)})}\bigr\|_{\mathrm{HS}}^2,
        \label{eq:level3}
      \end{equation}
      where $C_n = \kappa^{-2}\,|S_n|^{-2}$ depends only on the pipeline
      normalisation and the BFS geometry, and not on the pair $\{c, q-c\}$.

      In particular, $\sigma^{\mathrm{can}}_{\mathrm{pair}}(n)$ is the pullback of
      the canonical Hilbert--Schmidt quadratic form on $\mathrm{End}(V_\rho)$
      under $v_c^{(n)} \mapsto \Phi_{q,\rho}(v_c^{(n)})$, and therefore defines a
      rank-one trajectory in $\mathrm{End}(V_\rho)$ whose growth exponent coincides
      with $\delta_{\mathrm{pair}}$.
    \end{theorem}

    \begin{proof}
      Set $u^{(n)} := \Phi_{q,\rho}(v_c^{(n)})$ and
      $w^{(n)} := \Phi_{q,\rho}(v_{q-c}^{(n)})$.
      By condition~(a), $w^{(n)} = \overline{u^{(n)}}$.
      By condition~(b), $\|u^{(n)}\|^2 = d_\rho^{-1}\|v_c^{(n)}\|^2$ and
      $\|w^{(n)}\|^2 = d_\rho^{-1}\|v_{q-c}^{(n)}\|^2$.
      By~\eqref{eq:HS-rank1},
      \begin{equation}
        \|u^{(n)} \otimes \overline{w^{(n)}}\|_{\mathrm{HS}}^2
        = \|u^{(n)}\|^2 \cdot \|w^{(n)}\|^2
        = d_\rho^{-2}\,\|v_c^{(n)}\|^2 \cdot \|v_{q-c}^{(n)}\|^2.
      \end{equation}
      By~\eqref{eq:norm-sigma-relation}, in the pre-saturation regime,
      $\|v_c^{(n)}\|^2 \cdot \|v_{q-c}^{(n)}\|^2
      = \kappa^2\,|S_n|^2\,\sigma_c(n)\,\sigma_{q-c}(n) + O(\varepsilon_n)$.
      The factor $d_\rho^{-2}$ is absorbed by setting $C_n = d_\rho^2\,\kappa^{-2}|S_n|^{-2}$,
      which gives~\eqref{eq:level3}.
      The growth exponent of $C_n \|\cdot\|_{\mathrm{HS}}^2$ equals that of
      $\sigma^{\mathrm{can}}_{\mathrm{pair}}$ by the same argument as
      Proposition~\ref{sec:level1}.
    \end{proof}

    \begin{remark}
      The constant $C_n$ is uniform across pairs and does not affect the power-law
      exponent.
      The structural content of Theorem~\ref{thm:level3} is therefore entirely
      carried by the embedding $\Phi_{q,\rho}$.
      Hypothesis~\ref{hyp:embedding} remains open.
      O27~\cite{Beau2026a31} exhibits, under the spinor carrier supplied conditionally by O23, the
      canonical representation action $\rho \colon \mathfrak{su}(2) \to \mathrm{End}(V_\rho)$; the
      corresponding intertwiner space is one-dimensional, so an intertwiner is unique only up to a
      scalar.
      That construction does not furnish the vector lift $\mathcal{V}_q \to V_\rho$ required here.
      There is no non-zero $\mathfrak{su}(2)$-equivariant map from the adjoint to the fundamental,
      so an equivariant adjoint-factorising realisation is excluded (O27, Theorem~4.1(iv)); a
      quotient-factorising linear lift that is not equivariant is not excluded.
      Test~4 of §\ref{sec:tests} measures a property of the carrier $H_{\mathrm{eff}}$, not of $V_\rho$:
      O29~\cite{Beau2026a33} finds finite-data diagnostics compatible with the adjoint carrier
      $H_{\mathrm{eff}} \simeq \mathfrak{su}(2)$, without identifying that carrier, which is itself
      conditional (O27) on O23's supplied carrier and O27's admissible-saturation hypothesis; the spin-$\tfrac{1}{2}$
      sector itself is selected by minimality, not by
      Test~4.
    \end{remark}

  \subsection{Candidate sector}
    \label{sec:candidate}

    The admissibility hierarchy of~\cite{Beau2026a1} selects the spin-$\frac{1}{2}$
    sector of $2I$ as the minimal non-abelian admissible sector: dimension $d_\rho = 2$,
    Laplacian eigenvalue $\lambda_{1/2} = 18$ on the standard $2I$ Cayley graph.
    Under Hypothesis~\ref{hyp:embedding}, the proportionality constant in
    Theorem~\ref{thm:level3} is $d_\rho^{-2} = \tfrac{1}{4}$, and
    $\Phi_{q,\rho}(v_c^{(n)}) \otimes \overline{\Phi_{q,\rho}(v_{q-c}^{(n)})}$
    takes values in $\mathrm{End}(V_\rho) \simeq \mathbb{C}^{2 \times 2}$.
    Condition~(a) confines these values to a proper subspace.

    \begin{proposition}[Symmetric-square bound under involution equivariance]
      \label{prop:sym-bound}
      Assume condition~(a) of Hypothesis~\ref{hyp:embedding}, with the conjugation
      in~\eqref{eq:hyp-involution} taken in a fixed basis of $V_\rho$, and let
      $d = \dim_{\mathbb{C}} V_\rho$.
      Then for every shell $n$ the Level~III outer product is the symmetric square
      \[
        \Phi_{q,\rho}(v_c^{(n)}) \otimes \overline{\Phi_{q,\rho}(v_{q-c}^{(n)})} = u\,u^{\mathsf T},
        \qquad u := \Phi_{q,\rho}(v_c^{(n)}),
      \]
      so the trajectory spans a subspace of $\mathrm{Sym}^2(V_\rho)$, the complex symmetric
      $d \times d$ matrices, of dimension $d(d+1)/2$.
      The covariance
      $\mathcal{C}^{\Phi}_c := N^{-1}\sum_n \mathrm{vec}(\widehat{M}_n)\,\mathrm{vec}(\widehat{M}_n)^\dagger$
      of the embedded products
      $\widehat{M}_n := \Phi_{q,\rho}(v_c^{(n)}) \otimes \overline{\Phi_{q,\rho}(v_{q-c}^{(n)})}$,
      formed as in~\eqref{eq:covariance}, has rank $r^{\Phi}_{\mathrm{eff}} \leq d(d+1)/2$.
      For $d = 2$ this bound is $3$, and the rank-four target $d_\rho^2 = 4$ is excluded.
    \end{proposition}

    \begin{proof}
      By condition~(a), $w := \Phi_{q,\rho}(v_{q-c}^{(n)}) = \overline{u}$, hence $\overline{w} = u$.
      The matrix of $u \otimes \overline{w}$ has entries $u_a \overline{w_b} = u_a u_b$, which are
      symmetric in $a, b$; thus $u \otimes \overline{w} = u\,u^{\mathsf T} \in \mathrm{Sym}^2(V_\rho)$.
      The complex symmetric $d \times d$ matrices form a space of dimension $d(d+1)/2$.
      The range of $\mathcal{C}^{\Phi}_c$ is contained in the span of the vectors
      $\mathrm{vec}(\widehat{M}_n)$, hence in $\mathrm{vec}(\mathrm{Sym}^2(V_\rho))$.
      Since $d(d+1)/2 < d^2$ for $d \geq 2$, the value $d^2$ cannot be attained.
    \end{proof}

    Proposition~\ref{prop:sym-bound} is a statement about the contract of
    Hypothesis~\ref{hyp:embedding} as written.
    It does not refute an embedding without condition~(a), an embedding carrying a different
    anti-linear involution, or another construction of the pair trajectory, and it does not
    decide whether Hypothesis~\ref{hyp:embedding} holds.
    It fixes what a Level~III rank test can decide under that contract: the target $d_\rho^2 = 4$
    of the original Test~4 criterion (§\ref{sec:test4}) is not attainable there, for any data.
    For $d_\rho = 2$ the bound is $3$, the same value as the modal rank observed in
    $\mathrm{End}(H_{\mathrm{eff}})$: a rank equal to $3$ therefore cannot separate a Level~III trajectory
    from the adjoint reading of $H_{\mathrm{eff}}$, and it cannot be inverted to $d_\rho = 2$
    (O29~\cite{Beau2026a33}).

  \subsection{Structural interpretation}
    \label{sec:interpretation}

    At Level~I, $\sigma_{\mathrm{pair}}$ and the Hilbert--Schmidt norm share a growth
    exponent: the identification is asymptotic.
    At Level~II, $\sigma_{\mathrm{pair}}$ is the Hilbert--Schmidt norm up to a
    normalisation: the identification is structural modulo convention.
    At Level~III, $\sigma_{\mathrm{pair}}$ \emph{is} the Hilbert--Schmidt norm squared
    of a rank-one element of $\mathrm{End}(V_\rho)$, symmetric by
    Proposition~\ref{prop:sym-bound}, with the normalisation fixed by
    representation theory: this is a theorem under Hypothesis~\ref{hyp:embedding},
    not a conjecture.

  \section{Diagnostic criteria and their scope}
  \label{sec:tests}

  The identification proposed in §\ref{sec:conjecture} leads to computable diagnostics.
  Tests~1--3, and the form of Test~4 in the measured carrier, can be evaluated from the stored
  Gram--Schmidt data without modifying the O25 pipeline; the Level~III forms of Tests~4 and~5 require
  the embedding $\Phi_{q,\rho}$, which Hypothesis~\ref{hyp:embedding} leaves open.
  Tests~1 and~2 hold by construction and check the implementation.
  The conjugate-pair covariance of Test~4 computable from these data is formed in the measured
  carrier $H_{\mathrm{eff}}$ rather than in $\mathrm{End}(V_\rho)$
  (Remark~\ref{rem:criterion-reform}); under condition~(a) of Hypothesis~\ref{hyp:embedding}, the
  Level~III form of Test~4 is bounded by Proposition~\ref{prop:sym-bound}.

  \subsection{Test 1: Involution compatibility}
    \label{sec:test1}

    At the scalar level, the conjugation identity $\rho_{q-c} = \overline{\rho_c}$
    implies $\delta_{q-c} = \delta_c$~\cite{Beau2026a21}.
    At the matrix level, the dictionary maps $c \leftrightarrow q-c$ to
    $\widetilde{M}_n \leftrightarrow \widetilde{M}_n^\dagger$ identically, for any vectors
    (criterion~(i) of §\ref{sec:dict-map}).
    The Frobenius distance
    \begin{equation}
      \Delta_n(c) := \bigl\|\widetilde{M}_n(c)^\dagger
      - \widetilde{M}_n(q-c)\bigr\|_{\mathrm{HS}}
      \label{eq:test1}
    \end{equation}
    must vanish up to numerical precision for all $n \in [n_0, n_1]$.

    \begin{criterion}[Involution]
  $\Delta_n(c) < \varepsilon_{\mathrm{mach}}$ for all tested pairs and shells.
    \end{criterion}

    Since $\Delta_n(c)$ vanishes by construction, Test~1 is a consistency check on the
    implementation, like Test~2: a non-zero value signals a computational error and does not
    discriminate Levels~II and~III.
    The discriminating content of involution compatibility is condition~(a) of
    Hypothesis~\ref{hyp:embedding}, which concerns the embedded vectors and is not testable from
    the Weil-side vectors alone.

  \subsection{Test 2: Gauge invariance}
    \label{sec:test2}

    The factor $r(c,q)$ is a normalisation artefact of the pipeline~\cite{Beau2026a21}.
    The dictionary absorbs it by working with $\widetilde{M}_n$ rather than
    with individual normalised vectors.

    \begin{criterion}[Gauge invariance]
  Under the rescaling $v_c^{(n)} \mapsto \lambda\, v_c^{(n)}$,
  $v_{q-c}^{(n)} \mapsto \lambda^{-1} v_{q-c}^{(n)}$ for $\lambda > 0$,
  $\|\widetilde{M}_n\|_{\mathrm{HS}}^2$ is invariant.
    \end{criterion}

    This follows from~\eqref{eq:HS-rank1} by direct computation:
    $\|\lambda u \otimes \overline{\lambda^{-1} v}\|_{\mathrm{HS}}^2
    = \lambda^2 \|u\|^2 \cdot \lambda^{-2} \|v\|^2
    = \|u\|^2 \cdot \|v\|^2$.
    Test~2 is a consistency check on the construction, expected to hold identically.

  \subsection{Test 3: Exponent stability}
    \label{sec:test3}

    Level~I (Proposition~\ref{sec:level1}) predicts $\delta_{\mathrm{pair}} = \delta_{\mathrm{HS}}$.

    \begin{criterion}[Exponent stability]
  For each prime $q \in \{29, 61, 101, 151\}$, compute
  \begin{equation}
    \delta_{\mathrm{HS}}(q)
    = \mathrm{slope}\bigl(
    \log \|\widetilde{M}_n\|_{\mathrm{HS}}^2 \;\mathrm{vs}\; \log(n+1)
    \bigr)
  \end{equation}
  over the window $[n_0, n_1]$ used in O25~\cite{Beau2026a29} (Table~2 therein).
  The difference $|\delta_{\mathrm{HS}}(q) - \delta_{\mathrm{pair}}(q)|$
  must lie within the regression standard error reported in O25.
    \end{criterion}

    This test is the numerical anchor for the entire construction.

  \subsection{Test 4: Effective dimension of the trajectory}
    \label{sec:test4}

    This test is stated below in its original form; Proposition~\ref{prop:sym-bound} and
    Remark~\ref{rem:criterion-reform} fix what it can decide.
    For each conjugate pair $\{c, q-c\}$, construct the empirical covariance operator
    \begin{equation}
      \mathcal{C}_c = \frac{1}{N} \sum_{n=n_0}^{n_1}
      \mathrm{vec}(\widetilde{M}_n)\,\mathrm{vec}(\widetilde{M}_n)^\dagger,
      \quad N = n_1 - n_0 + 1,
      \label{eq:covariance}
    \end{equation}
    where $\mathrm{vec}$ flattens $\widetilde{M}_n$ into a column vector.

    \begin{criterion}[Effective dimension --- original formulation]
  Let $r_{\mathrm{eff}}$ be the numerical rank of $\mathcal{C}_c$
  (number of singular values above threshold
  $10^{-10} \cdot \sigma_1(\mathcal{C}_c)$).
  \begin{itemize}
    \item If $r_{\mathrm{eff}} = 4$ for all pairs: consistent with the
    spin-$\frac{1}{2}$ sector candidate ($d_\rho = 2$,
    $\dim\mathrm{End}(V_\rho) = 4$).
    \item If $r_{\mathrm{eff}} = d_\rho^2$ for some other fixed $d_\rho$:
    Level~III holds for a different irrep of $2I$; the candidate sector
    must be revised.
    \item If $r_{\mathrm{eff}}$ varies across pairs or primes: Level~III fails;
    the identification holds at most at Level~II.
  \end{itemize}
  Read at Level~III, the criterion concerns the covariance $\mathcal{C}^{\Phi}_c$ of the embedded products
  and its rank $r^{\Phi}_{\mathrm{eff}}$ (Proposition~\ref{prop:sym-bound}).
  Under condition~(a) of Hypothesis~\ref{hyp:embedding} and for $d_\rho \geq 2$, neither of the first two
  outcomes can occur for that rank: Proposition~\ref{prop:sym-bound} bounds $r^{\Phi}_{\mathrm{eff}}$ by
  $d_\rho(d_\rho+1)/2 < d_\rho^2$.
  This criterion is reinterpreted in Remark~\ref{rem:criterion-reform} below,
  following O29~\cite{Beau2026a33}: the conjugate-pair covariance is formed in $\mathrm{End}(H_{\mathrm{eff}})$
  rather than $\mathrm{End}(V_\rho)$, so no measured rank is a rank in $\mathrm{End}(V_\rho)$ or tests the
  target $d_\rho^2 = 4$, and the modal value
  $r_{\mathrm{eff}} = 3$ is finite-data evidence compatible with the supplied adjoint (spin-$1$) carrier
  $H_{\mathrm{eff}}$, not an identification of it.
    \end{criterion}

    \begin{remark}[What Test~4 measures --- O29]
      \label{rem:criterion-reform}
      The conjugation identity $\rho_{q-c} = \bar\rho_c$ (O17; recalled in O18 Theorem~3.1)
      does not by itself relate the projected coordinates of conjugate blocks.
      Under the additional projected-coordinate relation $\pi_{q-c}(v) = \overline{\pi_c(v)}$, which
      O29~\cite{Beau2026a33} tests and finds satisfied only approximately by the stored data, every
      outer product is complex symmetric.
      That algebraic restriction alone does not exclude rank $4$ inside the six-dimensional space
      $\mathrm{Sym}(H_{\mathrm{eff}}, \mathbb{C})$.
      In the tested samples the modal rank in $\mathrm{End}(H_{\mathrm{eff}})$ is three, while seed-sensitive
      single-sample exceptions reach ranks four to nine at $q = 101$ and five at $q = 211$; none of these
      ranks is a measurement in $\mathrm{End}(V_\rho)$.
      Breaking the per-vector locking restores rank four only in an auxiliary two-dimensional PCA
      truncation, at $q \in \{29, 61\}$, not in $\mathrm{End}(V_\rho)$.
      Under condition~(a) of Hypothesis~\ref{hyp:embedding}, by contrast, the rank-four target in
      $\mathrm{End}(V_\rho)$ is excluded for any data (Proposition~\ref{prop:sym-bound}); the O29
      samples, formed in $H_{\mathrm{eff}}$, do not bear on that bound.
      At the checkpoints examined in O29, $r_{\mathrm{eff}} = 3$ is the modal numerical value in
      $\mathrm{End}(H_{\mathrm{eff}})$, not an invariant; at $q \in \{29, 61\}$ the symmetric squares show a
      numerical $6 \to 3$
      Veronese collapse, and the fitted constraint forms have commutators of numerical rank three in a
      chosen $\mathfrak{so}(3)$ basis with a one-dimensional computed commutant.
      These diagnostics are compatible with the adjoint carrier $H_{\mathrm{eff}} \simeq \mathfrak{su}(2)$
      (spin-$1$, vector) conditional on O23's supplied carrier and O27's admissible-saturation
      hypothesis~\cite{Beau2026a31}; they neither identify
      that carrier from the data nor prove exact irreducibility.
      Test~4 does not identify the spin-$\tfrac{1}{2}$ sector $V_\rho$, which is not reachable
      equivariantly from the adjoint --- $\mathrm{Hom}_{\mathfrak{su}(2)}(\mathbf{3},\mathbf{2}) = 0$
      --- and requires an independent observable; the spin-$\tfrac{1}{2}$
      selection rests on the minimality argument of §\ref{sec:candidate} (the minimal non-abelian admissible
      sector), not on Test~4.
    \end{remark}

    \begin{remark}
      The matrix $\widetilde{M}_n$ is reconstructed from the stored vectors
      $v_c^{(n)}$ and $v_{q-c}^{(n)}$ by outer product.
      The covariance~\eqref{eq:covariance} and its SVD scale as $O(N d^4)$
      with $d$ the effective dimension; for $d = 2$ this is negligible.
      No additional BFS computation or Weil decomposition is required.
    \end{remark}

  \subsection{Test 5: Universality across pairs}
    \label{sec:test5}

    Level~III requires that all admissible pairs for a given $q$ probe the same
    representation sector.

    \begin{criterion}[Universality]
  The singular value spectra of $\mathcal{C}_c$ and $\mathcal{C}_{c'}$ must agree
  up to numerical fluctuations for all admissible pairs $c, c'$ at the same
  prime $q$.
    \end{criterion}

    Read at Level~III, the criterion compares the spectra of $\mathcal{C}^{\Phi}_c$ and
    $\mathcal{C}^{\Phi}_{c'}$ (Proposition~\ref{prop:sym-bound}).
    A pair-dependent spectrum would indicate that different pairs project onto
    different subspaces of $\mathrm{End}(V_\rho)$, which is incompatible with a
    canonical identification.
    The inter-pair concentration of $\delta_{\mathrm{pair}}$ established in
    O25 (coefficient of variation dropping from $6.1\%$ to $1.5\%$
    over $q \in \{29, 61, 101, 151\}$) provides indirect support for this test.

  \subsection{Decision table}
    \label{sec:decision}

    \begin{center}
      \begin{tabular}{p{0.52\textwidth}|p{0.40\textwidth}}
        \hline
        Outcome & Interpretation \\
        \hline
        Test~1 or Test~2 fails & Implementation error: both hold by construction \\
        Test~3 fails & Dictionary fails; §\ref{sec:dictionary} must be revised \\
        Tests~1--3 pass, $r^{\Phi}_{\mathrm{eff}}$ varies across pairs or primes & Level~III fails; at most
        Level~II (original criterion) \\
        Tests~1--3 pass, $r^{\Phi}_{\mathrm{eff}}$ fixed, Test~5 fails & At most Level~II; Level~III not
        supported \\
        Tests~1--3 and~5 pass, $r^{\Phi}_{\mathrm{eff}}$ fixed & Consistent with Level~III; under condition~(a) the
        fixed value is at most $d_\rho(d_\rho+1)/2$ ($3$ for $d_\rho = 2$), so it cannot equal $d_\rho^2$
        and does not by itself fix $d_\rho$ \\
        \hline
      \end{tabular}
    \end{center}

    \begin{remark}
      The rows are stated for the covariance $\mathcal{C}^{\Phi}_c$ of the embedded products and its rank
      $r^{\Phi}_{\mathrm{eff}}$, which require the embedding $\Phi_{q,\rho}$; Hypothesis~\ref{hyp:embedding}
      leaves that embedding open, so no row beyond the first two can be evaluated on current data.
      Tests~1 and~2 hold by construction and do not discriminate between the remaining rows.
      Under condition~(a) of Hypothesis~\ref{hyp:embedding}, no row can support Level~III through
      $r^{\Phi}_{\mathrm{eff}} = d_\rho^2$ (Proposition~\ref{prop:sym-bound}).

      The O28 and O29 data are formed in $\mathrm{End}(H_{\mathrm{eff}})$ and include no run of
      Tests~1--3, so they occupy no row of the table.
      O28~\cite{Beau2026a32} reported $r_{\mathrm{eff}} = 3$ across the pairs of its perimeter at
      $q \in \{61, 101, 151\}$.
      In the single-sample runs of O29~\cite{Beau2026a33}, rank three is modal but not invariant: it holds
      for $32/50$ pairs at $q = 101$, the others giving ranks four to nine, and for $103/105$ pairs at
      $q = 211$, the other two giving rank five; the multi-sample follow-ups give modal rank three for
      every exceptional pair.
      The modal rank therefore stays below four, the exceptions being seed-sensitive, and the
      locking-broken control restores rank four only in an auxiliary PCA truncation at $q \in \{29, 61\}$;
      none of these ranks is a measurement in $\mathrm{End}(V_\rho)$.
      These data are compatible with the adjoint (spin-$1$) carrier $H_{\mathrm{eff}} \simeq \mathfrak{su}(2)$,
      conditional on O23's supplied carrier and O27's admissible-saturation hypothesis, which they do not
      identify; a bare effective rank does not fix a representation dimension.
    \end{remark}

  \section{Implications}
  \label{sec:implications}

  \subsection{From bilinear observable to canonical norm}
    \label{sec:impl-norm}

    Under Hypothesis~\ref{hyp:embedding}, Theorem~\ref{thm:level3} establishes that
    $\sigma_{\mathrm{pair}}$ is the Hilbert--Schmidt norm squared of a rank-one element
    of $\mathrm{End}(V_\rho)$ selected by the admissibility thread
    $Q_8 \subset 2I \subset \mathrm{SU}(2)$, with normalisation fixed by $d_\rho$ alone.
    The emergence of the pair observable is a \emph{quadratic completion}: the passage
    from a linear object $v_c^{(n)} \in V_\rho$ to a rank-one operator
    $\Phi_{q,\rho}(v_c^{(n)}) \otimes \overline{\Phi_{q,\rho}(v_{q-c}^{(n)})}
    \in \mathrm{End}(V_\rho)$ endowed with its natural norm, which lies in $\mathrm{Sym}^2(V_\rho)$ by
    Proposition~\ref{prop:sym-bound}.

  \subsection{Boundary of the $\beta^*$ interpretation}
    \label{sec:impl-beta}

    O24~\cite{Beau2026a28} proves that vertical fibre structure does not perturb the
    observable rank entering $\delta_{\mathrm{pair}}$.  This fibre-side theorem remains intact.
    It does not establish the further capacity-to-rate prescription
    \begin{equation}
      \beta^* \approx \frac{1}{\delta_{\mathrm{pair}} + \tfrac{1}{2}}.
    \end{equation}
    That expression combines the changing-degree LPS growth equation with a capacity
    measured on the fixed-degree Heisenberg cascade; neither its square-root carrier nor
    its Born--Infeld factor is native to the latter substrate~\cite{Beau2026sgn}.
    Under Level~III, $\delta_{\mathrm{pair}}$ still characterises the decay of a
    Hilbert--Schmidt trajectory in $\mathrm{End}(V_\rho)$.  O27~\cite{Beau2026a31}
    constrains the admissible representation sector, but that conditional sector reading does
    not supply the missing growth equation.  Consequently the interval obtained by
    inserting measured values into the reciprocal is retained only as a phenomenological
    comparison; it has no derived representation-theoretic meaning for $\beta^*$.

  \subsection{Selection of the admissible sector}
    \label{sec:impl-sector}

    The spin-$\frac{1}{2}$ sector emerges as a candidate not from external physical
    input, but from internal admissibility constraints~\cite{Beau2026a1}.
    Test~4 of §\ref{sec:test4} yields finite-data evidence compatible with the adjoint carrier
    $H_{\mathrm{eff}} \simeq \mathfrak{su}(2)$ (spin-$1$, vector) conditional on O23's supplied carrier and O27's
    admissible-saturation hypothesis,
    without identifying it; the spin-$\tfrac{1}{2}$ selection rests on minimality and is not decided by
    Test~4 (O29~\cite{Beau2026a33}).
    The target space $\mathrm{End}(V_\rho) \simeq \mathbb{C}^{2 \times 2}$ is the ambient space of the Level~III
    products, whose trajectory lies in the three-dimensional $\mathrm{Sym}^2(V_\rho)$ under condition~(a)
    (Proposition~\ref{prop:sym-bound}); a direct test of $V_\rho$ requires an independent observable and a typed
    bridge to
    $V_\rho$, both left open in O29, and the dictionary of §\ref{sec:dictionary} remains valid throughout.

  \subsection{Role of non-injectivity}
    \label{sec:impl-noninj}

    The construction makes explicit how non-injectivity of the projection
    $\Pi : \chi \to \mathcal{O}$ induces a quadratic structure at the observable level.
    The pairing $\{c, q-c\}$ is not an arbitrary symmetrisation: it realises the
    parity involution, the candidate fibre structure of the pair sector
    (O18~\cite{Beau2026a22} establishes Born--Infeld parity as a covariance of the
    response family, with character $(-1)^{k}$ in the derivative order, and states the
    fibre identification as an open problem).
    Under that identification hypothesis, the Hermitian structure of
    $\mathrm{End}(V_\rho)$ is a projection-induced structure: it does not exist at
    the level of the substrate $\chi$, but emerges from the identification of
    conjugate configurations under $\Pi$.

  \subsection{Outlook}
    \label{sec:impl-outlook}

    The results of this note isolate a concrete bridge between the Weil
    representation of $\mathrm{Heis}_3(\mathbb{Z}/q\mathbb{Z})$, the admissibility
    thread $Q_8 \subset 2I \subset \mathrm{SU}(2)$, and the observable exponent
    $\delta_{\mathrm{pair}}$.
    Tests~1--3 of §\ref{sec:tests}, and the form of Test~4 in the measured carrier, can be executed on
    O25 data~\cite{Beau2026a29} without additional computation; the Level~III forms of Tests~4 and~5
    require the embedding $\Phi_{q,\rho}$.
    The effective dimension $r_{\mathrm{eff}} = 3$ was reported across the conjugate pairs of the
    perimeter of O28~\cite{Beau2026a32} at $q \in \{61, 101, 151\}$; the extended recomputation of
    O29~\cite{Beau2026a33} finds it modal rather than invariant ($32/50$ pairs at $q = 101$,
    $103/105$ at $q = 211$) and compatible with the adjoint (spin-$1$) carrier
    $H_{\mathrm{eff}} \simeq \mathfrak{su}(2)$ conditional on O23's supplied carrier and O27's admissible-saturation
    hypothesis, without identifying
    that carrier.
    Hypothesis~\ref{hyp:embedding} remains open.

    O27~\cite{Beau2026a31} exhibits the canonical representation action
    $\rho \colon \mathfrak{su}(2) \to \mathrm{End}(V_\rho)$ under the carrier supplied
    conditionally by O23.
    There is no non-zero $\mathfrak{su}(2)$-equivariant map from the adjoint
    representation to the fundamental representation, so an equivariant adjoint-factorising
    realisation of Hypothesis~\ref{hyp:embedding} is excluded (O27~\cite{Beau2026a31},
    Theorem~4.1(iv)).
    This does not exclude every quotient-factorising linear lift: the condition
    $\ker \Pi_{\mathrm{adm}} \subseteq \ker \Phi$ gives unique linear factorisation and no
    equivariance.
    Hypothesis~\ref{hyp:embedding} therefore remains open in that wider class.
    The implication chain
    \[
      \text{emergence} \Rightarrow \text{non-injectivity} \Rightarrow
      \text{pair structure} \Rightarrow \text{quadratic form} \Rightarrow \mathfrak{su}(2)
    \]
    has arrows of unequal type; its last step is the carrier supplied conditionally by O23, not a
    consequence of the three constraints.
    No carrier is identified numerically from O25 data: at $q \in \{29, 61\}$, O29~\cite{Beau2026a33} finds a
    numerical $6 \to 3$ Veronese collapse and finite-precision constraint diagnostics compatible with the supplied
    adjoint (spin-$1$) carrier $H_{\mathrm{eff}} \simeq \mathfrak{su}(2)$, and does not derive that carrier
    or prove exact irreducibility.
    The spin-$\tfrac{1}{2}$ sector $V_\rho$ is not identified by Test~4, and it is not reachable
    equivariantly from the adjoint.
    Under condition~(a) of Hypothesis~\ref{hyp:embedding}, any Level~III rank test is bounded by
    Proposition~\ref{prop:sym-bound}: the trajectory spans at most $\mathrm{Sym}^2(V_\rho)$, and the
    rank-four target is excluded under that contract.

    \begin{remark}[Test~4 in O29]
      \label{rem:test4-numerical}
      O29~\cite{Beau2026a33} reformulates Test~4 as a finite-data carrier diagnostic.
      The modal rank of the tested conjugate-pair samples in $\mathrm{End}(H_{\mathrm{eff}})$ is three, with
      seed-sensitive single-sample exceptions above it; no measured rank, modal or exceptional, is a rank in
      $\mathrm{End}(V_\rho)$ or tests the target $d_\rho^2 = 4$; the numerical $6 \to 3$
      Veronese collapse and $\mathfrak{so}(3)$ constraint pattern found at $q \in \{29, 61\}$ are compatible with the
      supplied adjoint
      (spin-$1$) carrier $H_{\mathrm{eff}} \simeq \mathfrak{su}(2)$ without identifying it (see
      Remark~\ref{rem:criterion-reform}).
      Numerically, $r_{\mathrm{eff}} = 3$ holds for every tested pair at $q \in \{29, 61, 151\}$
      (zero inter-pair variance), for $103/105$ pairs at $q = 211$, and as the modal value ($32/50$)
      at $q = 101$.
      The spin-$\tfrac{1}{2}$ sector $V_\rho$ is not decided by Test~4.
    \end{remark}

    \bibliographystyle{plain}
    \bibliography{cosmochrony-bibliography}

\end{document}
